\chapter{测试环境详细参数}\label{chap:detail-env}

本章给出\cref{chap:reinforce}中四项Fetch机械臂操作任务的完整参数规格。这些任务均在MuJoCo仿真引擎与Gymnasium-Robotics测试环境库上搭建，操作平台选用7自由度Fetch移动机械臂。如\cref{sec:rl-environment}所述，每项任务把输入和输出所有参数统一建模为动作空间、观测空间以及目标空间三个部分，并运用目标条件马尔可夫决策过程框架加以描述。以下各表按编号顺序列出每个维度的物理含义与对应库文档给定的取值范围。

Reach任务要求机械臂把末端执行器移动至空间中的指定目标点，不涉及对外部物体的操作。因此该任务的观测维度最为简单，仅包含10维机械臂本体运动学状态以及3维目标坐标，如\cref{tab:fetch-reach-full}所示。

Push、Slide以及PickAndPlace三项任务均要求机械臂与外部物体产生物理交互，因此观测空间在Reach的基础上加入了物体的空间位置、欧拉角、线速度、角速度，还有物体相对于末端执行器的相对位置与相对速度，构成25维的观测向量。三项任务的参数规格分别如\cref{tab:fetch-push-full}、\cref{tab:fetch-slide-full}以及\cref{tab:fetch-pickplace-full}所示。
\begin{table}[!htb]
  \caption{FetchReach-v4 任务参数规格表}\label{tab:fetch-reach-full}
  \centering
  \small
  \begin{tabular}{@{} llccll @{}} \toprule
    \textbf{类型} & \textbf{编号} & \textbf{参数名称} & \textbf{最小值} & \textbf{最大值} & \textbf{单位} \\ \midrule
    \textbf{动作} & 0-2 & 末端执行器在$(x, y, z)$方向的位移 & -1 & 1 & m \\
    & 3 & 夹爪控制(闭合不激活) & -1 & 1 & - \\ \midrule
    \textbf{观测} & 0-2 & 末端执行器在全局坐标系中的位置$(x, y, z)$ & -Inf & Inf & m \\
     & 3-4 & 左右夹爪手指的关节位移$(d_l, d_r)$ & -Inf & Inf & m \\
    & 5-7 & 末端执行器在三个方向的线速度$(v_x, v_y, v_z)$ & -Inf & Inf & m/s \\
    & 8-9 & 左右夹爪手指的线速度$(v_{l}, v_{r})$ & -Inf & Inf & m/s \\ \midrule
    \textbf{目标} & 0-2 & 末端执行器的目标坐标点$(x, y, z)$ & -Inf & Inf & m \\ \bottomrule
  \end{tabular}
\end{table}

\begin{table}[!htb]
  \caption{FetchPush-v4 任务参数规格表}\label{tab:fetch-push-full}
  \centering
  \small
  \begin{tabular}{@{} llccll @{}} \toprule
    \textbf{类型} & \textbf{编号} & \textbf{参数名称} & \textbf{最小值} & \textbf{最大值} & \textbf{单位} \\ \midrule
    \textbf{动作} & 0-2 & 末端执行器在$(x, y, z)$方向的位移 & -1 & 1 & m \\
    & 3 & 夹爪控制(闭合不激活) & -1 & 1 & - \\ \midrule
    \textbf{观测} & 0-2 & 末端执行器在全局坐标系中的位置$(x, y, z)$ & -Inf & Inf & m \\
    & 3-5 & 物体在全局坐标系中的位置$(x, y, z)$ & -Inf & Inf & m \\
    & 6-8 & 物体相对于末端执行器的相对位置$(x, y, z)$ & -Inf & Inf & m \\
    & 9-10 & 左右夹爪手指的关节位移$(d_l, d_r)$ & -Inf & Inf & - \\
    & 11-13 & 物体绕$(x, y, z)$三轴的旋转位姿 & -Inf & Inf & rad \\
    & 14-16 & 物体相对于末端执行器的线速度$(v_x, v_y, v_z)$ & -Inf & Inf & m/s \\
    & 17-19 & 物体角速度$(\omega_x, \omega_y, \omega_z)$ & -Inf & Inf & rad/s \\
    & 20-22 & 末端执行器线速度$(v_x, v_y, v_z)$ & -Inf & Inf & m/s \\
    & 23-24 & 左右夹爪手指速度$(v_{l}, v_{r})$  & -Inf & Inf & m/s \\ \midrule
    \textbf{目标} & 0-2 & 物体期望目标坐标$(x, y, z)$，$z$坐标固定 & -Inf & Inf & m \\ \bottomrule
  \end{tabular}
\end{table}


\begin{table}[!htb]
  \caption{FetchSlide-v4 任务参数规格表}\label{tab:fetch-slide-full}
  \centering
  \small
  \begin{tabular}{@{} llccll @{}} \toprule
    \textbf{类型} & \textbf{编号} & \textbf{参数名称} & \textbf{最小值} & \textbf{最大值} & \textbf{单位} \\ \midrule
    \textbf{动作} & 0-2 & 末端执行器在$(x, y, z)$方向的位移 & -1 & 1 & m \\
    & 3 & 夹爪控制(闭合不激活) & -1 & 1 & - \\ \midrule
    \textbf{观测} & 0-2 & 末端执行器的全局位置坐标$(x, y, z)$ & -Inf & Inf & m \\
    & 3-5 & 物体在全局坐标系中的位置$(x, y, z)$ & -Inf & Inf & m \\
    & 6-8 & 物体相对于末端执行器的相对位置$(x, y, z)$ & -Inf & Inf & m \\
    & 9-10 & 左右夹爪手指的关节位移$(d_l, d_r)$ & -Inf & Inf & - \\
    & 11-13 & 物体绕$(x, y, z)$三轴的旋转位姿 & -Inf & Inf & rad \\
    & 14-16 & 物体相对于末端执行器的线速度$(v_x, v_y, v_z)$ & -Inf & Inf & m/s \\
    & 17-19 & 物体角速度$(\omega_x, \omega_y, \omega_z)$ & -Inf & Inf & rad/s \\
    & 20-22 & 末端执行器线速度$(v_x, v_y, v_z)$ & -Inf & Inf & m/s \\
    & 23-24 & 左右夹爪手指速度$(v_{l}, v_{r})$  & -Inf & Inf & m/s \\ \midrule
    \textbf{目标} & 0-2 & 物体期望目标坐标$(x, y, z)$，$z$坐标固定 & -Inf & Inf & m \\ \bottomrule
  \end{tabular}
\end{table}

\begin{table}[!htb]
  \caption{FetchPickAndPlace-v4 任务参数规格表}\label{tab:fetch-pickplace-full}
  \centering
  \small
  \begin{tabular}{@{} llccll @{}} \toprule
    \textbf{类型} & \textbf{编号} & \textbf{参数名称} & \textbf{最小值} & \textbf{最大值} & \textbf{单位} \\ \midrule
    \textbf{动作} & 0-2 & 末端执行器在$(x, y, z)$方向的位移 & -1 & 1 & m \\
    & 3 & 夹爪控制(张开/闭合) & -1 & 1 & - \\ \midrule
    \textbf{观测} & 0-2 & 末端执行器的全局位置坐标$(x, y, z)$ & -Inf & Inf & m \\
    & 3-5 & 物体在全局坐标系中的位置$(x, y, z)$ & -Inf & Inf & m \\
    & 6-8 & 物体相对于末端执行器的相对位置$(x, y, z)$ & -Inf & Inf & m \\
    & 9-10 & 左右夹爪手指的关节位移$(d_l, d_r)$ & -Inf & Inf & - \\
    & 11-13 & 物体绕$(x, y, z)$三轴的旋转位姿 & -Inf & Inf & rad \\
    & 14-16 & 物体相对于末端执行器的线速度$(v_x, v_y, v_z)$ & -Inf & Inf & m/s \\
    & 17-19 & 物体角速度$(\omega_x, \omega_y, \omega_z)$ & -Inf & Inf & rad/s \\
    & 20-22 & 末端执行器线速度$(v_x, v_y, v_z)$ & -Inf & Inf & m/s \\
    & 23-24 & 左右夹爪手指速度$(v_{l}, v_{r})$  & -Inf & Inf & m/s \\ \midrule
    \textbf{目标} & 0-2 & 物体期望目标坐标$(x, y, z)$ & -Inf & Inf & m \\ \bottomrule
  \end{tabular}
\end{table}

上述四项任务在动作空间、奖励函数以及部分训练超参数上具备一定的共性，而在目标空间和动作空间约束上存在区别，下面从相同参数以及差异参数两个维度来进行对比。

相同参数包括四项任务的动作空间均为4维连续向量$\mathcal{A} = [-1, 1]^4$，前3个维度对应末端执行器在笛卡尔坐标系$(x, y, z)$方向上的增量位移，第4个维度对应夹爪开合控制。奖励函数均运用\cref{sec:reward-space}所述的稀疏奖励设计，距离容差阈值$\epsilon = 0.05$ m。

差异参数主要体现在三个方面。其一，目标空间的物理约束各不相同。Push与Slide任务中目标物体的$z$坐标固定于桌面高度，目标仅在水平面内分布；而PickAndPlace任务的目标坐标在三维空间中自由分布，涵盖桌面以上以及桌面以下的位置。其二，夹爪控制的激活条件有所区别。Reach、Push以及Slide任务中夹爪始终保持闭合状态，第4个动作维度不参与环境动力学计算；PickAndPlace任务中该维度被激活，负责控制两侧指尖的相对距离以实现抓取与释放。其三，各任务对与物体交互的要求与各自完成的复杂度不同。Reach仅涉及无接触直接机械臂末端控制；Push要求机械臂与物体持续接触推移；Slide则既包含了通过单次击打将物块打至远距离目标，也包含了持续推移使物块移动到近距离目标的两种行为模式；PickAndPlace则综合了抓取、搬运以及释放三个阶段。任务形式和复杂度上的这种差异直接影响了\cref{tab:fetch-hyperparams}中各任务的训练步数与网络架构选择。

\chapter{模型微调训练细节}\label{chap:finetune-details}

本章给出\cref{chap:finetune}中6个基座模型在单任务以及多任务微调设定下的训练损失曲线，用来补充\cref{sec:llm_experiment-results}中的性能分析。所有曲线均依据\cref{tab:training-config}所示的统一超参数配置记录，横轴为训练步数，纵轴为负对数似然损失值。下文涉及的具体损失数值均从训练曲线读图估计得出，需要由原始训练日志做进一步核实。

\section{训练损失曲线}
训练过程中采用的损失函数是标准的负对数似然损失，且仅针对辅助指令部分进行计算。\cref{fig:loss-single}展示了6个基座模型在四项独立任务以及多任务混合微调设定下的训练损失变化。各模型的初始损失均处于1.8\textasciitilde2.5的区间，反映出预训练语言模型对输出格式以及数值密集型动作序列的初始拟合能力较弱。随着训练的推进，所有曲线在前200步内出现了迅速下降，随后逐步趋于收敛。

\begin{figure}[!htb]
  \centering
  \begin{subfigure}[b]{0.31\textwidth}
    \centering
    \includegraphics[width=\textwidth]{figures/Reach}
    \caption{Reach}\label{fig:loss-reach}
  \end{subfigure}
  \hfill
  \begin{subfigure}[b]{0.31\textwidth}
    \centering
    \includegraphics[width=\textwidth]{figures/Push}
    \caption{Push}\label{fig:loss-push}
  \end{subfigure}
  \vspace{2pt}
  \begin{subfigure}[b]{0.31\textwidth}
    \centering
    \includegraphics[width=\textwidth]{figures/Slide}
    \caption{Slide}\label{fig:loss-slide}
  \end{subfigure}
  \hfill
  \begin{subfigure}[b]{0.31\textwidth}
    \centering
    \includegraphics[width=\textwidth]{figures/PickAndPlace}
    \caption{PickAndPlace}\label{fig:loss-pickplace}
  \end{subfigure}
  \begin{subfigure}[b]{0.31\textwidth}
    \centering
    \includegraphics[width=\textwidth]{figures/MultiTask}
    \caption{MultiTask}\label{fig:loss-multitask}
  \end{subfigure}
  \caption{六种基座模型在四项单任务微调设定下的训练损失曲线}\label{fig:loss-single}
\end{figure}

从\cref{fig:loss-single}的曲线形态中可以归纳出几方面的共性特征。在训练初期，所有模型在所有任务中均表现出迅速的损失下降，表明LoRA微调能够使模型快速捕捉输出格式。约1000步之后，各模型的损失逐步进入波动较小的平台期，下降速率明显减缓，该阶段模型主要在进一步优化对输出数值精度的学习。由于OLMo-2-1B和Phi-1.5每10步记录一次损失，而其余四个模型每50步记录一次，前两者的曲线在视觉上更为密集。收敛之后，损失在各模型间呈现出明显差异：Qwen3.5-0.8B在多数任务上收敛至最低或接近最低的损失，Qwen3-0.6B与Gemma3-1B紧随其后，三者均表现出较小的训练波动，表明其预训练表征与动作生成任务的适配性较好；Llama-3.2-1B与OLMo-2-1B收敛至中等水平的损失；Phi-1.5在多数任务上收敛至最高损失，训练过程中的波动幅度也最大，说明该模型的预训练目标与数值密集型动作生成任务之间存在较大差距。

从任务维度进一步观察，各模型在不同任务上的收敛表现同样存在差异。Reach任务仅涉及10维无接触运动学状态，观测维度最低、控制逻辑最为简单，各模型在该任务上初始损失范围为0.85\textasciitilde2.54，收敛速度较快，约4500步完成训练。其中Qwen3.5-0.8B收敛至0.63的最低损失，Qwen3-0.6B与Gemma3-1B紧随其后，均约为0.64，Llama-3.2-1B与OLMo-2-1B收敛至约0.81\textasciitilde0.86，Phi-1.5从最高初始损失2.54降至约0.98。Push任务的初始损失范围为0.35\textasciitilde2.05，普遍低于其余任务。OLMo-2-1B与Qwen3.5-0.8B收敛至约0.24\textasciitilde0.29的最低损失，Qwen3-0.6B与Gemma3-1B约为0.29，Phi-1.5初始损失最高达2.05，但相对降幅最大，最终收敛至约0.31，Llama3.2-1B收敛至约0.38。Slide任务由于任务复杂度较高，损失范围在所有任务中最大，初始损失0.88\textasciitilde2.52，最终损失0.69\textasciitilde1.13。Gemma3-1B收敛至0.69，Qwen3.5-0.8B为0.71，Qwen3-0.6B为0.73，Llama3.2-1B、OLMo2-1B与Phi-1.5分别收敛至0.91、0.99与1.13。PickAndPlace任务涉及三维空间中物体的抓取与搬运，各模型在训练初期约1000步内即从2.1以上迅速降至1.0以下，其中Gemma3-1B收敛至0.41，Qwen3.5-0.8B为0.43，Qwen3-0.6B为0.44，Llama3.2-1B为0.55，OLMo2-1B为0.60，Phi-1.5为0.72。

\section{多任务与单任务微调对比分析}

\cref{fig:loss-multitask}展示了多任务混合微调设定下的训练损失曲线。在该设定下，来自四项任务的训练样本被均匀混合后共同输入模型，模型需要在单次训练过程中同时学习多项任务的观测动作映射关系。与\cref{fig:loss-single}所示的单任务曲线相比，多任务训练损失曲线的整体趋势仍遵循迅速下降到逐步收敛的两阶段模式，但在几个方面呈现出差异。

在该设定下，训练扩展至15100步，来自四项任务的训练样本被均匀混合后共同输入模型，模型需要在单次训练过程中同时学习多项任务的观测—动作映射关系。与单任务设定相比，所有模型的收敛过程更为平缓。Qwen3.5-0.8B与Qwen3-0.6B取得最低的最终损失0.50，Gemma3-1B紧随其后，损失0.51，Llama-3.2-1B、OLMo-2-1B与Phi-1.5分别收敛至0.65、0.67与0.79。尽管训练数据来自混合的任务分布，所有模型均表现出单调的损失下降趋势，未出现任务间的灾难性干扰。

将\cref{fig:loss-single}与\cref{fig:loss-multitask}的观察结果对照\cref{tab:experiment-results}中的闭环控制成功率可知，收敛损失较低、训练波动较小的模型在实际控制任务中通常取得更高的成功率，该趋势在Reach、PickAndPlace等任务上体现较为一致；而收敛损失较高且波动剧烈的模型在闭环评估中的表现也相应较弱。需要指出的是，训练损失仅反映模型对训练集分布的拟合程度，与闭环控制性能之间仅具有参考性的相关关系，不能直接等价，但可作为模型选型的辅助参考依据。